\documentclass[a4paper,twoside]{article}

\usepackage{morewrites}

\usepackage[svgnames,dvipsnames,x11names]{xcolor}
\usepackage{graphicx}
\usepackage[english]{babel}
\usepackage[colorlinks=true, linkcolor=NavyBlue, urlcolor=NavyBlue, citecolor=NavyBlue]{hyperref}
\usepackage[a4paper]{geometry}
\usepackage{ts-fonts}
\usepackage{csquotes}
\usepackage{float}
\usepackage{graphicx}
\usepackage{seqsplit}

\usepackage{lastpage}
\usepackage{refcount}
\makeatletter
\def\ps@plain{
\let\@oddhead\@empty
\let\@evenhead\@empty
\def\@oddfoot{
\hfil
\ifnum\getpagerefnumber{LastPage}>1\relax
\thepage
\fi
\hfil}
\let\@evenfoot\@oddfoot
}
\makeatother

\title{Booby}
\author{Yves Chevallier}\date{November 16, 2025}
\begin{document}
\maketitle
\section{Introduction}
Boobies are seabirds in the genus \emph{Sula}, family Sulidae. They are
large, long-winged birds that plunge-dive for fish. The name \enquote{booby}
originates from the Spanish word \enquote{bobo}, meaning \enquote{stupid} or \enquote{clown},
due to the birds' apparent lack of fear of humans.
\begin{figure}[H]
\centering
\includegraphics[width=0.3\linewidth]{booby.png}
\caption{Booby}
\end{figure}
\section{Particularities}
Boobies have several distinctive features:
\begin{itemize}
\item They have brightly colored feet, which they use in mating displays.
\item They are known for their spectacular diving ability, plunging into
the water from great heights to catch fish.
\end{itemize}
\end{document}
